% SSSUP-FONTE-STORICA: fonti-locali/originali/2014-2020.zip
% Anno 2015, esercizio 2. Trascrizione redazionale dal PDF della prova scritta.
% Verificare editorialmente la trasposizione prima di considerarla definitiva.
\renewcommand{\Anno}{2015}
\renewcommand{\NumeroProblema}{2}

\renewcommand{\Testo}{
La Scuola Superiore Sant'Anna deve cablare un'aula in modo tale che:

\begin{quote}
L'impianto elettrico realizzato \textbf{garantisca} la presenza di \textbf{almeno una presa elettrica} (rappresentata dal simbolo in figura sotto) ad una distanza di \textbf{al più un metro} da \textbf{ogni punto} della stanza.
\end{quote}

L'impianto deve essere realizzato utilizzando una sola \textbf{presa fissa}, (disponibile gratuitamente e da collocare sul pavimento dell'aula ove ritenuto più opportuno), e un numero arbitrario di \emph{prese triple} e \emph{prolunghe}, schematizzate nella figura sottostante.

\begin{center}
% SSSUP-FIGURA-RICOSTRUITA: simboli elettrici da Scritto-MAT-2015.pdf.
\psset{unit=0.70cm}
\begin{pspicture}(-0.3,-1.1)(12.7,1.35)
  \psset{linewidth=.75pt}
  \psline(1.5,-.05)(1.5,.48)
  \psline(1.5,.48)(1.27,.74)
  \psline(1.5,.48)(1.73,.74)
  \psline(1.17,-.05)(1.83,-.05)
  \rput(1.5,-.67){presa fissa}
  \psline(5.2,-.05)(5.2,.27)
  \psline(5.2,.27)(4.89,.46)
  \psline(5.2,.27)(5.51,.46)
  \psline(5.2,.27)(5.2,.59)
  \psline(5.2,.59)(5.02,.76)
  \psline(5.2,.59)(5.38,.76)
  \psline(4.89,.46)(4.69,.68)
  \psline(4.89,.46)(4.65,.47)
  \psline(5.51,.46)(5.71,.68)
  \psline(5.51,.46)(5.75,.47)
  \psline(4.96,-.05)(5.44,-.05)
  \rput(5.2,-.67){presa tripla}
  \psline(8.4,.45)(10.0,.45)
  \psline[linestyle=dashed,dash=2pt 3pt](10.0,.45)(11.25,.45)
  \psline(11.25,.45)(12.3,.45)
  \psline(8.4,.45)(8.63,.69)
  \psline(8.4,.45)(8.63,.21)
  \psline(12.3,.45)(12.05,.68)
  \psline(12.3,.45)(12.05,.22)
  \rput(10.4,-.67){prolunga}
\end{pspicture}
\end{center}

Considerando tutte le prese come puntiformi, determinare:
\begin{enumerate}
\item[(a)] l'impianto di costo minimo che garantisca la copertura di una stanza di $2\times2$ metri nel caso in cui ogni presa tripla costi 1 Euro e una prolunga lunga $\ell\in\mathbb R$ metri costi $\ell$ Euro (si possono usare prolunghe di lunghezza arbitraria);

\item[(b)] l'impianto di costo minimo che garantisca la copertura di una stanza di $2\times2$ metri nel caso in cui ogni presa tripla costi 1 Euro e si abbia a disposizione \textbf{unicamente} prolunghe di un metro di lunghezza, ciascuna di costo 1 Euro;

\item[(c)] il lato $D$ della stanza quadrata di dimensione massima che è possibile coprire con prese triple (al costo di 1 Euro a pezzo) e prolunghe di lunghezza 1 metro (al costo di 1 Euro a pezzo), avendo a disposizione 6 Euro.
\end{enumerate}

\textbf{Commento.} In assenza di una soluzione esatta, sono accettabili soluzioni approssimate, a patto che la copertura sia garantita. In questo caso, sarà valutata positivamente anche la capacità di fornire una misura di \emph{distanza} dall'ottimo. Esempio: ``la soluzione XYZ costa 10 Euro, ma garantisco che non esiste soluzione che costi meno di 3 Euro.
}

\renewcommand{\Risposta}{}
\renewcommand{\Aiutino}{}
\renewcommand{\SoluzioneUfficiale}{}
\renewcommand{\SoluzioneRedazionale}{}
\renewcommand{\SoluzioniAlternative}{}
\renewcommand{\Osservazioni}{}
\renewcommand{\Commenti}{}
